\begingroup
\renewcommand{\arraystretch}{1.45}
\Needspace{40\baselineskip}
\InterpolationArticle{12}
Es lohnt sich, eine Transformation, die derjenigen analog ist, die wir in Artikel 4 erläutert haben,
auch auf die Formeln von Artikel 10 und dem vorangehenden Artikel anzuwenden, wenn dort $\alpha^m\cos\tfrac12s+\beta^m\sin\tfrac12s=0$
gesetzt wird; dadurch lassen sich diese beiden verschiedenen Fälle unter einer einzigen Formel zusammenfassen.
Die Anzahl der Größen $a,b,c,d\mathinner{.\,.\,.}$ wachse daher schrittweise
von einer $a$ auf zwei $a,b$, dann auf drei $a,b,c$ usw.; man entwickle abwechselnd
nach den Artikeln 10 und 11 den Wert, den $t$ unter jeder dieser Voraussetzungen annimmt, und ziehe schließlich
jeden Wert vom unmittelbar folgenden ab, und zwar so:
\[
\begin{aligned}
&A\\[5pt]
&\frac{\sin\tfrac12(t-b)}{\sin\tfrac12(a-b)}A\cos\tfrac12(t-a)\\[5pt]
&+\frac{\sin\tfrac12(t-a)}{\sin\tfrac12(b-a)}B\cos\tfrac12(t-b)\\[5pt]
&\frac{\sin\tfrac12(t-b)\sin\tfrac12(t-c)}{\sin\tfrac12(a-b)\sin\tfrac12(a-c)}A\\[5pt]
&+\frac{\sin\tfrac12(t-a)\sin\tfrac12(t-c)}{\sin\tfrac12(b-a)\sin\tfrac12(b-c)}B\\[5pt]
&+\frac{\sin\tfrac12(t-a)\sin\tfrac12(t-b)}{\sin\tfrac12(c-a)\sin\tfrac12(c-b)}C\\[5pt]
&\text{usw.}
\end{aligned}
\]
Hier ist die Differenz zwischen dem ersten und zweiten Wert
\[
=\sin\tfrac12(t-a)\cos\tfrac12(t-b)
\left\{\frac{A}{\sin\tfrac12(a-b)}+\frac{B}{\sin\tfrac12(b-a)}\right\}
\]
Die Differenz zwischen dem zweiten und dritten Wert ist
\[
\begin{aligned}
={}&\sin\tfrac12(t-a)\sin\tfrac12(t-b)\left\{
\frac{A\cos\tfrac12(a-c)}{\sin\tfrac12(a-b)\sin\tfrac12(a-c)}\right.\\[5pt]&\qquad\left.
+\frac{B\cos\tfrac12(b-c)}{\sin\tfrac12(b-a)\sin\tfrac12(b-c)}\right.\\[5pt]&\qquad\left.
+\frac{C}{\sin\tfrac12(c-a)\sin\tfrac12(c-b)}\right\}
\end{aligned}
\]
Setzt man diese Operation weiter fort, so ergibt sich leicht das allgemeine Gesetz. Setzt man nämlich
\SourceLeaf{294}
\[
\begin{aligned}
A'={}&\frac{A}{\sin\tfrac12(a-b)}+\frac{B}{\sin\tfrac12(b-a)}\\[5pt]
A''={}&\frac{A\cos\tfrac12(a-c)}{\sin\tfrac12(a-b)\sin\tfrac12(a-c)}\\[5pt]
&+\frac{B\cos\tfrac12(b-c)}{\sin\tfrac12(b-a)\sin\tfrac12(b-c)}\\[5pt]
&+\frac{C}{\sin\tfrac12(c-a)\sin\tfrac12(c-b)}\\[5pt]
A'''={}&\frac{A}{\sin\tfrac12(a-b)\sin\tfrac12(a-c)\sin\tfrac12(a-d)}\\[5pt]
&+\frac{B}{\sin\tfrac12(b-a)\sin\tfrac12(b-c)\sin\tfrac12(b-d)}\\[5pt]
&+\frac{C}{\sin\tfrac12(c-a)\sin\tfrac12(c-b)\sin\tfrac12(c-d)}\\[5pt]
&+\frac{D}{\sin\tfrac12(d-a)\sin\tfrac12(d-b)\sin\tfrac12(d-c)}\\[5pt]
\end{aligned}\]
\[\begin{aligned}
A''''={}&\frac{A\cos\tfrac12(a-e)}{\sin\tfrac12(a-b)\sin\tfrac12(a-c)\sin\tfrac12(a-d)\sin\tfrac12(a-e)}\\[5pt]
&+\frac{B\cos\tfrac12(b-e)}{\sin\tfrac12(b-a)\sin\tfrac12(b-c)\sin\tfrac12(b-d)\sin\tfrac12(b-e)}\\[5pt]
&+\frac{C\cos\tfrac12(c-e)}{\sin\tfrac12(c-a)\sin\tfrac12(c-b)\sin\tfrac12(c-d)\sin\tfrac12(c-e)}\\[5pt]
&+\frac{D\cos\tfrac12(d-e)}{\sin\tfrac12(d-a)\sin\tfrac12(d-b)\sin\tfrac12(d-c)\sin\tfrac12(d-e)}\\[5pt]
&+\frac{E}{\sin\tfrac12(e-a)\sin\tfrac12(e-b)\sin\tfrac12(e-c)\sin\tfrac12(e-d)}
\end{aligned}
\]
usw., so gilt
\[
\begin{aligned}
T={}&A+A'\sin\tfrac12(t-a)\cos\tfrac12(t-b)\\[5pt]
&+A''\sin\tfrac12(t-a)\sin\tfrac12(t-b)\\[5pt]
&+A'''\sin\tfrac12(t-a)\sin\tfrac12(t-b)\sin\tfrac12(t-c)\cos\tfrac12(t-d)\\[5pt]
&+A''''\sin\tfrac12(t-a)\sin\tfrac12(t-b)\sin\tfrac12(t-c)\sin\tfrac12(t-d)\\[5pt]
&+A^{\mathrm{V}}\sin\tfrac12(t-a)\sin\tfrac12(t-b)\sin\tfrac12(t-c)\\[5pt]&\qquad \sin\tfrac12(t-d)\sin\tfrac12(t-e)\cos\tfrac12(t-f)\\[5pt]
&+A^{\mathrm{VI}}\sin\tfrac12(t-a)\sin\tfrac12(t-b)\sin\tfrac12(t-c)\\[5pt]&\qquad \sin\tfrac12(t-d)\sin\tfrac12(t-e)\sin\tfrac12(t-f)\\[5pt]
&+\text{usw.}
\end{aligned}
\]
Diese Reihe ist auf so viele Glieder fortzusetzen, wie Werte der Funktion $X$
gegeben sind.

Die Koeffizienten $A',A'',A''',A''''$ usw. lassen sich auch mit dem folgenden Algorithmus berechnen.
Mit $B',B'',B''',B''''$ usw. werde bezeichnet, was aus ihnen jeweils wird, wenn
\[
a,\quad A,\quad b,\quad B,\quad c,\quad C,\quad d,\quad D\quad\text{usw.}
\]
der Reihe nach übergehen in
\[
b,\quad B,\quad c,\quad C,\quad d,\quad D,\quad e,\quad E\quad\text{usw.}
\]
Ferner sollen die Übergänge
\[\begin{gathered}
A',A'',A''',A''''\text{ usw.}\quad\text{in}\quad C',C'',C''',C''''\text{ usw.}\quad
\\[5pt]\text{oder in}\quad D',D'',D''',D''''\text{ usw. usw.}
\end{gathered}\]
erfolgen, wenn
\SourceLeaf{295}
\[
a,\ A,\ b,\ B,\ c,\ C,\ d,\ D\quad\text{usw.}
\]
der Reihe nach übergehen in
\[
c,\ C,\ d,\ D,\ e,\ E,\ f,\ F\quad\text{usw.}
\]
oder in
\[
d,\ D,\ e,\ E,\ f,\ F,\ g,\ G\quad\text{usw.}
\]
\hspace*{2em}usw.

Dann gilt
\[\begin{gathered}
A'=\frac{A-B}{\sin\tfrac12(a-b)},\\[5pt]
B'=\frac{B-C}{\sin\tfrac12(b-c)},\\[5pt]
C'=\frac{C-D}{\sin\tfrac12(c-d)}\quad\text{usw.}
\end{gathered}\]
\[\begin{gathered}
A''=\frac{A'\cos\tfrac12(a-c)-B'}{\sin\tfrac12(a-c)},\\[5pt]
B''=\frac{B'\cos\tfrac12(b-d)-C'}{\sin\tfrac12(b-d)}\quad\text{usw.}
\end{gathered}\]
\[\begin{gathered}
A'''=\frac{A''\cos\tfrac12(a-c)+A'\sin\tfrac12(a-c)-B''}{\sin\tfrac12(a-d)},\\[5pt]
B'''=\frac{B''\cos\tfrac12(b-d)+B'\sin\tfrac12(b-d)-C''}{\sin\tfrac12(b-e)}\quad\text{usw.}
\end{gathered}\]
\[\begin{gathered}
A''''=\frac{A'''\cos\tfrac12(a-e)-B'''}{\sin\tfrac12(a-e)},\\[5pt]
B''''=\frac{B'''\cos\tfrac12(b-f)-C'''}{\sin\tfrac12(b-f)}\quad\text{usw.}
\end{gathered}\]
\[\begin{gathered}
A^{\mathrm{V}}=\frac{A''''\cos\tfrac12(a-e)+A'''\sin\tfrac12(a-e)-B''''}{\sin\tfrac12(a-f)},\\[5pt]
B^{\mathrm{V}}=\frac{B''''\cos\tfrac12(b-f)+B'''\sin\tfrac12(b-f)-C''''}{\sin\tfrac12(b-g)}\quad\text{usw.}
\end{gathered}\]
\[
A^{\mathrm{VI}}=\frac{A^{\mathrm{V}}\cos\tfrac12(a-g)-B^{\mathrm{V}}}{\sin\tfrac12(a-g)},\qquad
B^{\mathrm{VI}}=\frac{B^{\mathrm{V}}\cos\tfrac12(b-h)-C^{\mathrm{V}}}{\sin\tfrac12(a-h)}\quad\text{usw.}
\]
usw.

Das Bildungsgesetz ist hier hinreichend ersichtlich, wenn man nur beachtet, dass die Zähler in den Werten für $A''$, $A'''$, $A''''$, $A^{\mathrm{V}}$, $A^{\mathrm{VI}}$ usw. (der Wert für $A'$ ist von dieser Regel auszunehmen) abwechselnd aus zwei oder drei Teilen bestehen.

\endgroup
